% =====================================================================
%  Draft Results text (framing). Numbers from appendix_results.tex
%  (re-evaluated checkpoints, <model>_eval.csv, 30 seeds everywhere).
%  All errors are MAE in raw units.
%  CBM sentences assume the rerun with per-concept MLPs keeps the same
%  ordering -- check before submission.
% =====================================================================

\subsection{Summary of findings}

We evaluate three claims: (i) a multiplicative gate $c_i\,\varphi_i(z)$ fails
when a concept is $0$, whereas anchor interpolation does not; (ii) on continuous
concepts, interpolation is more accurate and far cheaper than discretising the
concept into bins; and (iii) on real data, these benefits come at no cost in
predictive performance.

\paragraph{The zero failure of multiplicative embeddings.}
The zero test isolates the samples in which a concept equals $0$
(Tables~\ref{tab:app-mnist-zero} and~\ref{tab:app-pendulum-zero}). Because
groups can differ in difficulty, we compare each model's error on the zero
group with its own error on the non-zero group. For CEM (ours) and every binned
model the zero group is \emph{not} harder under full intervention
(task MAE, ArithmeticMNIST: $1.14$ vs.\ $1.34$; Pendulum: $0.102$ vs.\
$0.117$). For the multiplicative variants it is: on ArithmeticMNIST the
intervened error rises from $2.47$ to $3.26$ ($+32\%$) for CEM-Linear and from
$3.17$ to $4.47$ ($+41\%$) for CEM-Linear (no LN), and on Pendulum from $0.298$
to $0.385$ ($+29\%$) for CEM-Linear (no LN). Correcting \emph{only} the
zero-valued concept shows the same pattern: it reduces the error of CEM (ours)
from $2.61$ to $2.06$, but barely helps CEM-Linear ($3.27 \to 3.10$) and
\emph{increases} the error of CEM-Linear (no LN) (ArithmeticMNIST
$4.24 \to 5.06$; Pendulum $0.285 \to 0.344$): telling these models the correct
value $0$ erases the information carried by that concept's embedding.

The generative setting makes the mechanism visible
(Table~\ref{tab:app-zero-gen}). Object shape is not a concept, so it reaches the
decoder only through the concept embeddings. When the concepts of real images
are set to $0$, CEM keeps the shape of every object ($\geq 0.99$ shape accuracy
for any number of zeroed concepts), whereas CEM-Linear loses it progressively
($0.97$, $0.89$, $0.70$ for 2, 3, 4 zeroed concepts) and, with all five at $0$,
maps every input to the same image (shape accuracy $0.25$, i.e.\ chance).
A single zeroed concept is compensated by the remaining blocks in both settings
(Table~\ref{tab:app-orient}); the failure grows with the share of the
bottleneck that is zero. This explains why it is pronounced with two concepts
(ArithmeticMNIST, Pendulum) and not visible on LIDC-IDRI, where nine concepts
are rarely zero together.

\paragraph{Interpolation versus discretisation.}
Binning removes the zero problem as well, since no embedding is multiplied by
the concept value. Its cost is resolution and size. The number of parameters
of a binned bottleneck grows linearly with the number of bins $k$, since every
bin of every concept has its own embedding, whereas ours uses two anchors per
concept regardless of the resolution required (Table~\ref{tab:app-params}):
$1.5\times$, $2.5\times$ and $3.4\times$ the parameters of our model for
$k = 5, 10, 15$ on ArithmeticMNIST and Pendulum, and $1.7\times$, $2.8\times$
and $3.9\times$ on LIDC-IDRI. In addition, $k$ is a dataset-dependent hyperparameter, and a
binned model can only represent, predict or receive as an intervention the
$k$ bin centres.

These costs pay off only when the concepts are themselves discrete. The
ArithmeticMNIST concepts are digits, i.e.\ nine categories, and with $k \geq 10$
bins every digit gets its own embedding; here binning is the correct inductive
bias and the binned models are the most accurate (task MAE $1.55$ vs.\ $2.59$
for $k = 15$; intervened $0.84$ vs.\ $1.30$). With $k = 5$, which merges
digits, binning is still better without intervention ($2.17$ vs.\ $2.59$) and
tied with ours under intervention ($1.33$ vs.\ $1.30$, $p = 0.7$). On Pendulum, whose concepts are genuinely continuous
angles, our model is the most accurate on every metric: binned models have
$1.2$--$1.3\times$ its task error ($0.152$--$0.164$ vs.\ $0.126$) and
$1.3$--$1.4\times$ its intervention error ($0.154$--$0.166$ vs.\ $0.116$), and
their concept error is bounded below by the bin width, whatever the number of
bins we tried: the pendulum angle is off by $10.5$--$23.4$ (in the units of the
dataset) against $6.8$ for ours. On LIDC-IDRI, binning is
worse than every continuous CEM under intervention, and increasingly so with
more bins: balanced accuracy $0.857$, $0.846$ and $0.830$ for $k = 5, 10, 15$,
against $0.868$ for ours (all $p < 0.05$), with $1.7\times$, $2.8\times$ and
$3.9\times$ the parameters, and a higher concept error ($0.117$--$0.127$ vs.\
$0.115$ relative absolute error, all $p < 0.01$). The raw-unit errors per concept
(Table~\ref{tab:app-lidc-raw}) qualify this: binned models with $k \ge 10$ are
more accurate on texture, diameter and surface area, and worse on the other
ratings and on volume. Without intervention all models are tied. The Shapes3D results agree: the binned model reaches $0.933$
concept accuracy, against $0.9997$ for ours, on concepts whose true values lie
on a fixed grid.

\paragraph{Real data.}
On LIDC-IDRI the radiologists' concepts are subjective, noisy and do not
determine malignancy, and here the continuous CEM variants are statistically
indistinguishable under intervention (balanced accuracy $0.868$, $0.869$,
$0.872$ for CEM, CEM-Linear and CEM-Linear (no LN); differences to ours
$p \ge 0.14$). All three improve over CBM ($0.847$) by $2.1$--$2.5$ points
(ours vs.\ CBM: $p < 10^{-4}$). \todo{Check against the CBM
rerun.} Without intervention, ours is $1.2$ points below CEM-Linear (no LN)
($0.831$ vs.\ $0.843$, $p = 0.006$) and tied with CEM-Linear. The zero failure
does not appear on this dataset (Table~\ref{tab:app-lidc-zero}). We read this as
evidence that the representation we propose does not cost performance on real
data, while it removes a failure mode that the multiplicative variants exhibit
whenever zero-valued concepts carry a large share of the bottleneck.

\paragraph{Following interventions that contradict the input.}
Full intervention with the \emph{true} concepts favours the embedding models,
but these concepts agree with the input. The swap test
(Tables~\ref{tab:app-pendulum-swap} and~\ref{tab:app-lidc-swap}) instead feeds
each input the concepts of another sample. On Pendulum, where the target is a
function of the concepts, CBM follows them exactly (its head sees nothing
else), whereas every embedding model partly predicts from the input: the error
with respect to the new target (task MAE) is $0.44$ for CBM, $1.41$ for
CEM-Linear (no LN), $2.22$--$2.59$ for the binned models, $2.47$ for
CEM-Linear and $2.97$ for ours, and our predictions stay close to the
\emph{original} target (MAE $0.62$). On
LIDC-IDRI the pattern is weaker but the same (balanced accuracy against the new
label $0.66$--$0.67$ for the continuous CEMs, $0.70$--$0.73$ binned, $0.85$
CBM). Anchor interpolation therefore gives the most accurate interventions when
the intervened concepts are consistent with the input, and the least faithful
ones when they are not: the input-dependent anchors carry enough information
about the input to override a contradictory concept value. Counterfactual
training addresses exactly this case (Section~\ref{app:shapes3d}).
\todo{Decide how to present this; consider the swap loss for the benchmark.}
